\documentclass[SE,lsstdraft,authoryear,toc]{lsstdoc}
\input{meta}

% Package imports go here.

\usepackage{subcaption}
\usepackage{hyperref}

% Local commands go here.

%If you want glossaries
%\input{aglossary.tex}
%\makeglossaries

\title{Determining the Accuracy and Precision of Astrometric Positions and Covariances from the LSST Science Pipelines and the Vera C. Rubin Observatory Using LSST's Data Preview 1}

% This can write metadata into the PDF.
% Update keywords and author information as necessary.
\hypersetup{
    pdftitle={Determining the Accuracy and Precision of Astrometric Positions and Covariances from the LSST Science Pipelines and the Vera C. Rubin Observatory Using LSST's Data Preview 1},
    pdfauthor={Tom J. Wilson},
    pdfkeywords={}
}

\graphicspath{{./Plots/}}

% Optional subtitle
% \setDocSubtitle{A subtitle}

\input{authors}

\setDocRef{SITCOMTN-180}
\setDocUpstreamLocation{\url{https://github.com/lsst-sitcom/sitcomtn-180}}

\date{\vcsDate}

% Optional: name of the document's curator
% \setDocCurator{The Curator of this Document}

\setDocAbstract{%
This work extends SITCOMTN-159 as an investigation into the agreement between astrometric positions and their corresponding covariances as produced by the LSST Science Pipelines.
Here we use the Data Previews as the earliest pre-cursor to the main LSST survey.
These will serve as the initial benchmark of the accuracy and precision of the LSST astrometry, and guide users of the data on where the data are well aligned and what corrections may need to be made.

The analysis compares the pipeline-derived precisions, as described by the covariance reported in the Data Preview, to the variation in measured position of each dataset entry.
The measured position variation is computed by comparing LSST positions to two proxies for the ground truth, both high-quality astrometric datasets: \textit{Gaia} and the \textit{Hubble} Source Catalog.
Fitting these distributions, for a set of similar detections -- all with the same pipeline covariance, brightness, all in regions of the sky with the same density of background detections and same LSST depth, and so on -- we can build a set of empirical uncertainties to compare to those generated by the data reduction process.
In the case of both the \texttt{Source} and \texttt{Object} tables in DP1 we see a consistent under-estimation of the precisions as compared with experimental position-measurement scatter.
Some of the under-estimation of the precisions can be explained by an anomalous scaling of precisions with \texttt{Object} table coadd depth, suggesting systematic factors affecting the reported location of detections that isn't captured by pure photon-noise statistics.
}

% Change history defined here.
% Order: oldest first.
% Fields: VERSION, DATE, DESCRIPTION, OWNER NAME.
% See LPM-51 for version number policy.
\setDocChangeRecord{%
  \addtohist{0.1}{2026-03-04}{Initial draft of DP1 results.}{Wilson}
}


\begin{document}

% Create the title page.
\maketitle
% Frequently for a technote we do not want a title page  uncomment this to remove the title page and changelog.
% use \mkshorttitle to remove the extra pages

% ADD CONTENT HERE
% You can also use the \input command to include several content files.

\section{Introduction}
The two primary pieces of information extracted from astronomical images are the position and brightness of the objects in the region of sky observed.
The fundamental task of the image extraction pipeline is to turn the two-dimensional array of pixel counts into a database of unique detections and their parameters.
Along with removing non-astrophysical noise from the astrophysical signal, it must quantify both the ``best'' determinations of these parameters and the certainty with which we made that determination.
In the particular case of astrometry, this encapsulates the belief we have in the best-fit or highest-likelihood measurement being the true location of the source on the sky.

This higher-order information is necessary for a large swathe of astronomical problems, such as the problem of constructing object motions, but more fundamentally the precisions with which object locations were determined are crucial for the combining of two or more such catalogues together.
While the counterpart assignment problem can be solved in limited cases with just best-estimate positions of objects, modern algorithms for tackling the problem of determining the counterparts to catalogue objects (e.g. \citealp{Budavari:2008aa}; \citealp{2018MNRAS.473.5570W}; \citealp{Pineau:2017aa}) use detection uncertainties to complement the quoted position.
In doing so, likelihoods of the same object being detected twice, at the detections' respective positions with their respective detection certainties, can be compared to the chance encounter between two unrelated objects, neither of which is observed in the opposing telescope image.

It is this meta-science case that we are particularly concerned with here.
It is necessary for any astrometric measurements made by a data reduction pipeline to be in alignment with the precisions with which those measurements were made for those data to be combined with another dataset.
The catalogues generated must contain both accurate positions and reliable uncertainty quantification for the composite, catalogue-merged data that go into any and all scientific investigations to be robust; the conclusions drawn from mis-specified data will themselves be questionable.

To that end, in this work we extend previous investigations into the astrometry of the Vera C. Rubin Observatory's Legacy Survey of Space and Time (LSST; \citealp{Ivezic2019}), performed on pre-cursor simulated data \citep{sitcomtn-159}.
We utilise the first Data Preview from the Observatory ahead of the main survey, DP1, to characterise the accuracy and precision of its astrometry.

\section{Ensemble Astrometry and Astrometric Precisions}
As discussed by \citet{sitcomtn-159}, determining the agreement, or lack thereof, between the measured position of any astrophysical detection and the covariance matrix parameterising the belief in the measurement is generally not possible.
As is often the case in astronomy, we only have the single estimate of any given source at any given time, and do not know where the object really is to quantify how well-measured the position is -- if we did, the question of cross-matching and counterpart assignment would be a lot simpler!
So, in general, we must consider ensembles of objects, all with the same (or similar enough) parameters -- signal-to-noise ratio (SNR), flux, crowding levels, and so on -- and make the assumption that these objects all obey the same underlying statistical distribution in those parameters.
In particular, we must ensure that any systematic effects that may influence the measured positions of our objects are accounted for \citep[e.g.,][]{2017MNRAS.468.2517W,2018MNRAS.481.2148W,Wilson2023}, such that the precisions we are considering are purely those of the centroid from the pipeline.
While these additional effects, such as proper motion, influence both the measured position and, through hyper-parameter uncertainty, the uncertainty of the measured position, we are chiefly concerned here with measuring the effect the data reduction pipeline has on the catalogue.
We must therefore in essence subtract these other effects from our measurements to make the comparison as cleanly as possible.

Hence, if we have a set of $N$ detections, all with the same nuisance parameters and all with the same believed astrometric precision, we should be able to use these objects as a collective to test the pipeline-derived uncertainty.
If we also have some measure of each object's \textit{error} -- its deviation from its true position -- then we expect the distribution of these errors to follow a particular functional form.
The number of object errors in some small box $(\delta x,\ \delta y)$ at some deviation $(\Delta x,\ \Delta y)$ away from their ``real'' position is expected to follow a Gaussian distribution (see left-hand panel, Figure \ref{fig:gauss_v_ray}).
The parameter that controls this distribution is its standard deviation, $\sigma$; this can be computed in each dimension as e.g. $\sigma^2 = \frac{1}{N} \sum_{i=1}^N(x_i - \mu)^2$ with $\mu = \frac{1}{N}\sum_{i=1}^Nx_i$.
An example of this distribution is shown in the right-hand panel of Figure \ref{fig:gauss_v_ray}.
Alternatively, if we wish to circularise our problem and consider simply all objects some \textit{radial} extent away from zero deviation, using e.g. $\Delta r = \sqrt{\Delta x^2 + \Delta y^2}$, in some small annulus $\delta r$, then we instead find our set of errors described by a Rayleigh distribution (also shown in the right-hand panel, Figure \ref{fig:gauss_v_ray}).
Crucially in either case we use the same $\sigma$ consideration\footnote{Although the more general covariance matrix accounts for non-circular uncertainties, if the correlation between right ascension and declination is low enough, and the precisions in both sky axes roughly equal, we can use a single value to describe the scatter. This will have further advantage in future work in crowded fields in particular, where the ``extra,'' non-centroid components of the Astrometric Uncertainty Function \citep[AUF;][]{2017MNRAS.468.2517W} are fully circularised.}, and so in what follows either one- or two-dimensional parameterisation could be used.

\begin{figure}
  \centering
  \includegraphics[width=0.95\columnwidth]{{gaussianrayleigh}.pdf}
  \caption{Schematic showing the difference between the Gaussian and Rayleigh distributions.
           The left-hand figure shows the typical scatter in measured positions away from some ``true'' (unknown) position, in both sky axes.
           The right-hand figure shows the mathematical representation of this density, in terms of number of objects per unit area, in the black line, as a Gaussian distribution.
           Equivalently we can represent this as number of objects per unit distance from the origin, which transforms the function into a Rayleigh distribution, shown in the red line.
           In essence, as we map from a small sky square to a circular annulus inside which we determine our object counts, the Rayleigh distribution decreases towards the origin due to the decreasing area being considered.
           Figure reproduced from \citet{sitcomtn-159}.}
  \label{fig:gauss_v_ray}
\end{figure}

Following \citet{sitcomtn-159}, we wish to consider this simple question: do the standard deviations of the position errors agree with the pipeline-derived uncertainties?

\section{LSST's Data Preview 1, \textit{Gaia}, and the \textit{Hubble} Source Catalog}
Unfortunately, \citet{sitcomtn-159} had an easier time of it than usual in computing astrometric errors -- the simulated dataset considered had, by definition, a truth table.
For real data we will have to fall back on the next best thing: a significantly more astrometrically precise dataset; in this case, as in so many cases, \textit{Gaia} \citep{Collaboration2021}.
However, as \textit{Gaia} is a relatively bright survey, at least compared with LSST, we shall also make use of the \textit{Hubble} Source Catalog (HSC; \citealp{Whitmore2016}) to probe the dynamic range below that which \textit{Gaia} can detect objects.

We will use \textit{Gaia} and the HSC as our best-effort ``truth'' tables from which we can derive LSST errors in the following sections.

\subsection{Data Preview 1}
This first LSST data product contains seven pointings that were taken with the LSST commissioning camera, LSSTComCam.
These are spread across the southern sky, and represent a few key science areas the LSST will tackle; see RTN-011 \citep{Guy2025} for details, particularly tables 5 and 6.
As we wish to maximise our number statistics and ensure the astrometric parameterisations are valid across the entire range of science areas, we shall use all of these in our investigation.

We extract positions (\texttt{coord\_ra} and \texttt{coord\_dec}) and precisions (\texttt{coord\_raErr}, \texttt{coord\_decErr}, and \texttt{coord\_ra\_dec\_Cov}), computing the geometric mean of the corresponding semi-major and semi-minor axes of the covariance\footnote{This preserves the area of the uncertainty ellipse.}.
\texttt{coord\_raErr} and \texttt{coord\_ra\_dec\_Cov} were corrected for an issue in the DP1 dataset where the $\cos{\delta}$ correction parameter was over-applied to the right ascension precisions\footnote{Due to the spherical coordinate system, one must account for the declination of the observation when handling the right ascension component of the precision. Essentially, the 360-degree loop around right ascension gets smaller the further towards the poles the observation takes place. However, this correction was applied twice to the DP1 dataset erroneously; the issue is documented \href{https://dp1.lsst.io/products/known_issues_and_subtleties.html\#incorrect-geometric-scaling-in-dp1-right-ascension-errors}{here}, at least for the \texttt{coord\_raErr} component. The covariance can be given as $\sigma_{xy} = \rho \sigma_x \sigma_y$ with $\rho$ the correlation between $x$ and $y$ uncertainties, which should be in the range [-1, 1]. Investigations revealed a $\cos{\delta}$ term is introduced into the correlation parameter if \texttt{coord\_ra\_dec\_Cov} is not corrected for this effect when \texttt{coord\_raErr} is, although this is not reported on the known-issues page of DP1.}.
We additionally desire magnitudes and SNRs, which are computed through \texttt{[A]\_[B]Flux} and \texttt{[A]\_[B]FluxErr}, with \texttt{[A]} any of the six filters (u, g, r, i, z, or y).
\texttt{[B]} is either \texttt{cModel} or \texttt{psf}, depending on whether we wish to tune for extended or point sources respectively, depending on the Galactic latitude of the pointing.
Magnitudes are derived through \texttt{astropy} functionality, as recommended\footnote{Via \texttt{mag = (flux * u.nJy).to(u.ABmag)}.}, and SNR calculated as the ratio of flux to flux ``err'' (uncertainty, rather than error in the sense we are using it in this work).

Entries were filtered for basic flags, being removed from the final dataset if the astrometric uncertainty was \texttt{NaN} or negative.
Detections were also removed if all filters -- where again \texttt{[A]} represents any one of the six filters -- had SNR or magnitude uncertainty that was \texttt{NaN} or negative.
Additionally detections were removed if any one of the flags \texttt{[A]\_flag}, \texttt{[A]\_pixelflags\_edge}, \texttt{[A]\_pixelflags\_offimage}, or \texttt{[A]\_pixelflags\_saturated} were set in all bands.

\subsection{\textit{Gaia}}
We extracted the \textit{Gaia} positions (\texttt{ra} and \texttt{dec}) from the DR3 catalogue.
We similarly extracted the covariance matrix for each detection (\texttt{ra\_error}, \texttt{dec\_error}, and \texttt{ra\_dec\_corr}; again precisions rather than measurement errors), and circularised the uncertainties similar to DP1.
Magnitudes (\texttt{phot\_*\_mean\_mag} for \texttt{g}, \texttt{bp}, and \texttt{rp}) and SNRs (\texttt{phot\_*\_mean\_flux\_over\_error}) were also extracted.
We also used the proper motions in right ascension and declination from the database, keeping \texttt{pmra}, \texttt{pmdec}, and \texttt{ref\_epoch} (although this is fixed to \texttt{2016.0} for all sources by definition of the \textit{Gaia} motion solutions).
The coordinates of each \textit{Gaia} object were propagated to the main epoch of the LSST observations, as stored in \texttt{[A]\_epoch}, approximately \texttt{2024.9}.

Objects were removed from the dataset if the \texttt{ruwe} metric was greater than five, if all magnitudes were \texttt{NaN}, or -- in particular, as an extra criterion -- if the BP filter was fainter than 20.3 magnitudes (i.e., the object is removed if $G$ is \texttt{NaN} and $RP$ is \texttt{NaN} and $BP$ is \texttt{either} \texttt{NaN} or too faint).

\subsection{\textit{Hubble} Source Catalog}
In addition to the \textit{Gaia} dataset, we will be using the HSC as an extra, fainter dataset of comparison sources to compute LSST errors from.
The HSC is a meta-catalogue, containing a merger of the numerous disparate \textit{Hubble} observations over its long history.
We therefore do not have the usual level of quality flags with which to ensure a clean dataset; we will simply have to rely on the high quality of the observations and our false-positive rejection methods.

We will always require the sky coordinates of each unique observation, \texttt{matchra} and \texttt{matchdec}.
However, there are literally dozens of different filters across all HSC observations, the unique filters used in at least one \textit{Hubble} observation over its 30 years of observing.
We chose to extract the six most commonly used, for simplicity, which cover more than 90\% of detections in the dataset: ACS's \texttt{F475W}, \texttt{F555W}, \texttt{F606W} and \texttt{F814W} filters (with prefix \texttt{a\_}), and WFC3's \texttt{F110W} and \texttt{F160W} IR filters (prefix \texttt{w3\_}).
We additionally extract the photometric uncertainty -- formally the Median Absolute Deviation (MAD) of repeat measurements -- as e.g. \texttt{a\_f475w\_mad}, converting to SNR as needed.
Given the long baseline of observations, we also record the epoch of each set of observations via the \texttt{StopMJD} column.

\section{Validating the Astrometric Positions and Precisions of DP1}
\label{sec:validating-positions-precisions}
In this section we follow the same procedure as \citet{sitcomtn-159} used to test the LSST Science Pipelines against Operational Rehearsal 3 data.
Hence, the first step of the process is to generate a cross-match between our test data (LSST's DP1) and our truth table (\textit{Gaia} or HSC).
We wish to avoid, ironically, biasing our results by using any ``complicated'' cross-match algorithms that may require astrometric precisions to give us a sample of errors to compare to those very precisions.
We therefore use the only algorithm we can in such a situation: the nearest-neighbour match.

Each LSST object has its on-sky separation to each e.g. \textit{Gaia} detection calculated, and the one with the smallest distance is recorded as its nearest neighbour.
However, of course, the entire point of the field of probabilistic cross-matching algorithms is that the nearest-neighbour match is very sensitive to false positive interlopers in crowded fields, so we mitigate this by requiring that the \textit{Gaia} detection also has the LSST source as its closest neighbour, symmetrically.
This may decrease the completeness of our sample in such cases, but here we only really require a pure sample, so the trade-off is a good one.
For the cases where we consider DP1's time-series \texttt{Source} table, we will match each unique visit ID to the truth table separately and record the ``error'' of each separately.
The maximum radius of nearest neighbour separation was set to five arcseconds, but this is well above the lengthscales of separation we see and so we do not have to worry about this particular parameter and whether it influenced the results.

The experiment from here is quite straightforward.
For each pointing (see \citealp{Guy2025}, table 5, for details), we generate a series of histograms of distributions of reliable position errors, one-dimensional on-sky distances.
We then fit this distribution of errors with a Rayleigh distribution, representing the statistical model for the centroid precisions\footnote{For deeper or more crowded observations we will combine this centroid precision component of the full AUF with others -- chief of which is the perturbation due to unresolved contaminant detections in crowded fields \citep{2018MNRAS.481.2148W}. For now, however, we find that these DP1 data do not suffer any of these effects significantly and, similar to \citet{sitcomtn-159}, we will not consider them further in our DP1 investigation.}.
We also include in our fit a false-positive model.
These separations are generally seen on much broader lengthscales than our LSST position errors typical present -- being related to the inverse-square root of the catalogue sky densities -- and we include them for completeness to remove the small number of nuisance false matches we may expect to see at very faint magnitudes, but do not believe their exact treatment affects the results here and do not discuss them further.
Ultimately, for each pointing, we have a series of ``input'' uncertainties -- the pipeline-derived precisions -- and a series of ``output'' uncertainties -- the scatter-based error standard deviations -- from which we can begin to assess the agreement between the measured positions and the quoted certainty we have in those measurements.

\subsection{The \texttt{Object} Table Error-Precision Relation}
\label{sec:object-err-precision}
The main conclusion \citet{sitcomtn-159} drew from analysis of simulated LSST data was a breakdown in the relationship between best-fit position errors and pipeline-derived precisions for the deep-stack, coadded \texttt{Object} table, so it is there that we begin.
The \texttt{Object} table represents the detections from a pipeline run on the deepest possible set of observations, adding flux from all single-visit exposures together where they overlap in the sky to maximise SNR.
This is then the static sky component of the LSST data release, a complement to the time-series capabilities the observing cadence will bring.

Our main investigation will proceed with \textit{Gaia} as our truth table.
This dataset is the most precise astrometric reference we have, and it is all-sky, so maximises our overlap with our limited DP1 observations.
The significantly more precise astrometric positions of \textit{Gaia} allow us to effectively ignore one side of the equation in the distribution of our counterpart errors, and hence much more simply consider these errors as being ascribed entirely to LSST\footnote{As \citet{Ivezic2014} show, eventually LSST is expected to equal \textit{Gaia}'s precision by end-of-WFD-survey depth, but clearly this will not be true for the Data Previews or even its Early Science program, and we can use this to our advantage while we might!}.
However, it is relatively bright, so we will only be able to investigate the top end of the LSST dynamic range.

Following the method described above, for each pointing we generate a series of data sub-sets, at different brightnesses.
These magnitude slices run from 14th magnitude through to 24th magnitude, in increments of 0.1 magnitudes.
Given the number of exposures in each pointing in DP1 (\citealp{Guy2025}, table 6), we selected the $r$ band in particular for this experiment.
24th magnitude is a few magnitudes below the \textit{Gaia} completeness limit, but we wanted to ensure we had as many valid data points as possible; in practice the expected \textit{Gaia} completeness limit between 21st and 22nd magnitude and the roughly zero \textit{Gaia}-LSST colour meant we didn't see many cross-catalogue counterpart assignments at fluxes this faint.
Additionally, at least for such small DP1 pointings, we did not anticipate sufficient number statistics to draw meaningful conclusions for the very brightest of these sources -- especially given LSST's saturation limit -- but again kept the optimistic dynamic range just in case, with an eye to the much larger sky coverage of Data Preview 2.
For each magnitude we select all detections within 0.05 magnitudes (i.e., $r\pm0.05$ for $r$ in \{14, 14.1, 14.2, 14.3, ...\}).
A cut was also made on pipeline-derived astrometric precision to ensure we removed any outlier precisions that may indicate poor quality detections that would likely also come with position error outliers.
Finally, we also required a good number of detections of each magnitude-astrometric precision slice used, to give a robust least-squares fit.

\begin{figure}
  \centering
  \includegraphics[width=0.95\columnwidth]{{auf_fits_G_95-25_object_gaia_r18.9}.pdf}
  \caption{Comparison between astrometric precision model and distribution of measured-truth position separations.
           Black errorbars show the distribution of radial separations between \textit{Gaia} and LSST \texttt{Object} detections of $r=18.9$ for the LSST DP1 ``Rubin SV 95 -25'' pointing.
           The black line shows the expected distribution of separations based on quoted LSST astrometric precisions, while red and cyan lines show different best-fitting distributions.
           In this case the quoted astrometric uncertainties, as derived by the pipeline, are too small (black vs red/cyan solid lines).
           $F$ is the fraction of false positives in the best-fit model in each case, while $H$ is a parameter related to the crowding element of the uncertainty function, unimportant in these low-density regions.}
  \label{fig:95-25_obj_g_18.9}
\end{figure}

An example fit is shown in Figure \ref{fig:95-25_obj_g_18.9}, for the case of the DP1 ``Rubin SV 95 -25'' pointing (\citealp{Guy2025}, table 5), for objects within 0.05 magnitudes of $r=18.9$ with symmetric nearest neighbours in the \textit{Gaia} catalogue.
The \textit{Gaia}-DP1 errors are plotted in the black errorbars, showing Poissonian counting statistics uncertainties for each bin.
Notice the underestimation of the LSST astrometric precisions relative to the lengthscale of the measured \textit{Gaia}-DP1 errors.
The quoted centroid uncertainty of these LSST data is 0.4 milli-arcseconds (plotted in the poorly fitting solid black line), while the scatter in the centroids suggests a more reasonable \~3 milli-arcseconds (mas) centroid precision (in which the cyan solid line is a least-squares fit to the data).

As mentioned above, we require an additional truth table to probe those magnitudes below \textit{Gaia}'s $G \gtrsim 22$, and so here we turn to the HSC.
Very similar to the steps we use for \textit{Gaia}, we create DP1-HSC cross-matches, then step from 14th to 28th magnitude -- slightly below where we now expect \textit{LSST}'s completeness limit to be -- in increments of 0.1 magnitudes, taking $r\pm0.05$ cuts as before.
We filter for outlier pipeline precisions and make a cut on slices with insufficient number of matches, once again.

\begin{figure}
  \centering
  \includegraphics[width=0.95\columnwidth]{{auf_fits_F_ExCDFS_object_hubble_r24.4}.pdf}
  \caption{Comparison between astrometric precision model and distribution of measured-truth position separations.
           All symbols and lines have the same meaning as Figure \ref{fig:95-25_obj_g_18.9}, except the matches are between the \textit{Hubble} Source Catalog and DP1's \texttt{Object} table, and LSST was filtered on $r=24.4$. 
           The pointing in this case is the Extended Chandra Deep Field South.
           Note the small second peak at approximately 0.2 arcseconds, most likely un-accounted for proper motion.}
  \label{fig:excdfs_obj_h_24.4}
\end{figure}

We show an example of these results in Figure \ref{fig:excdfs_obj_h_24.4}.
Here the data are from the ``Extended Chandra Deep Field South'' pointing, with the brightness cut centered around $r=24.4$.
We still see a similar story to Figure \ref{fig:95-25_obj_g_18.9}, just with differing numbers: a pipeline-derived precision of 26 mas compares with an error-based scatter of roughly 64 mas.
The pipeline-derived covariance is once again plotted in the black solid line against the best-fit Rayleigh distribution (cyan line).

There are two complications with the \textit{Hubble} data that do not affect the \textit{Gaia} data that must be considered, however.
First, where the \textit{Gaia} data are negligible in the centroid error (and hence empirical precision) budget, the HSC data are not quite so astrometrically robust.
We find (see Appendix \ref{sec:hsc_appendix} for the details) that the HSC astrometric precisions are of order 15 mas, which must be subtracted in quadrature\footnote{The measurement error is a sum of the distances two measurements are away from zero, the sum of two vectors, and the mathematical form of the sum of two distributions is a convolution. The convolution of two Gaussians is analytic, itself a Gaussian with covariance matrix the sum of the two covariance matrices describing the respective distributions. In the simple case of circular uncertainties this reduces to the sum-in-quadrature of the scale parameters, and hence we subtract-in-quadrature to find the value we are interested in.}, giving an LSST-error of $\sqrt{64^2 - 15^2} \approx 62$ mas\footnote{At such faint magnitudes the assumed-constant HSC precisions don't contribute very much to the error budget, but this is important for brighter LSST objects and so must be taken into account.}.
Second, we are unable to account for the time baseline between the sets of observations, as we are using \textit{Gaia}'s proper motion information.
This is apparent in Figure \ref{fig:excdfs_obj_h_24.4} in the small peak in the data at approximately 0.2 arcseconds, showing a cluster of objects with similar on-sky motion, slightly larger than our scatter.
These objects not having their separations minimised serves to slightly increase the lengthscales of our errors, slightly inflating the empirically measured precision validations.
However, as can be seen, the number of such high-motion -- or, at least, high enough to be measurable -- objects is small, and the proper motion of objects decreases with physical distance from the Earth and hence, on average, brightness as observed from Earth.
Given that the \textit{Hubble} data are most useful in this faint regime, starting at around 23rd magnitude, this bias is going to be small at worst and become negligible in the most valuable cases.

\begin{figure}
  \centering
  \includegraphics[width=0.95\columnwidth]{{hubble_corrected_lsst_errors_object}.pdf}
  \caption{Summary of pipeline-derived uncertainties and experimental error-based scatter for DP1 astrometry.
           Each colour-symbol combination represents a different pointing in DP1, fit to a different truth table.
           The plusses use \textit{Gaia} as a truth table, and hence are a relatively bright and precise dataset with good overlap in all DP1 observations, while the crosses use \textit{Hubble} data to investigate LSST uncertainty, with only two pointings of common observing area that probe much fainter.
           The HSC data have had the non-negligible \textit{Hubble} astrometric precision removed while we assume the \textit{Gaia} precisions effectively do not contribute to the nearest neighbour distances.
           The left-hand panel shows the data in a linear scale while the right-hand panel shows the same datapoints on a logarithmic scale, to highlight the very brightest observations.}
  \label{fig:object_summary}
\end{figure}

Figure \ref{fig:object_summary} shows the full results of the \texttt{Object} table analysis.
Here we see -- in linear space on the left-hand side, dominated by the \textit{Hubble} pointings, and on a log-log scaling on the right, highlighting the bright \textit{Gaia} results -- the full extent of the investigation.
Unlike \citet{sitcomtn-159}, we do not see a power-law slope relationship between precision and error; in particular, we do not see precisions that are ever larger than the experimental error, as was reported for the OR3 dataset.
As \citeauthor{sitcomtn-159} had no way to investigate this scaling relation in the first place and could not explain the result in any case, we can do nothing here except chalk the original dependency up to an artefact of the simulation and consider the results we see here.

In the bright, DP1-\textit{Gaia} analysis, we see linear scalings of precision and error -- something of the much simpler form $y = m x + n$, or $y^2 = (m x)^2 + n^2$ -- for these data.
The only exception is the orange datapoints of Figure \ref{fig:object_summary}, from the ``95 -25'' pointing.
This set of observations suffers by far the most from scatter in the SNR of an object at a given magnitude, and this likely explains the ``wall'' of constant pipeline-quoted astrometric precision; this is investigated further in Section \ref{sec:object-snr-error-precision}.

With sufficient number statistics in many bins, we can further reduce these data into a parameterised model, as \citet{sitcomtn-159} did for the OR3 \texttt{Source} table.
The fits here aren't quite as tight, but ultimately for the above model of $y^2 = (mx)^2 + n^2$, in which a linear slope-corrected precision is added in quadrature with a missing systematic to explain the experimental errors, we find $n$ in the range 2-10 mas, but a linear slope $m$ in the range 1-5.
These hyper-parameter fits are shown in Figures \ref{fig:95-25_obj_g_18.9} and \ref{fig:excdfs_obj_h_24.4}, red lines, for reference.
The main conclusions here are limited -- largely by sky coverage, but also by \textit{Gaia} dynamic range -- but ultimately the conclusion is that the covariances are under-estimated as compared with the errors.

For the HSC, we have even less sky overlap, barely a single pointing, but we can probe to much fainter magnitudes.
Here we continue to see, for most of the LSST brightness range, an under-estimation of the covariances relative to the empirical scatter standard deviations.
Given both the lack of validated astrometric uncertainties for the \textit{Hubble} data and the lack of repeat pointings to validate the relationship we see, we are hesitant to draw much of a conclusion from these results.
A sum-in-quadrature hyper-parameter fit hints at a near-unity slope but implies $n \simeq 50\,\mathrm{mas}$, much higher than the hypothetical LSST systematics seen as compared to \textit{Gaia}.
Some or all of this could be attributed to the HSC errors, although, as discussed in Appendix \ref{sec:hsc_appendix}, the bright HSC data do not show scatter as compared with \textit{Gaia} data of even half of this gulf.
So, we cannot so easily explain away this large systematic uncertainty.
However, we have not accounted for any proper motion between HSC and LSST while we have done so in the HSC-\textit{Gaia} check, and in our main ExCDFS pointing there is anywhere from seven to 22 years epoch difference.
Thus, depending on the fraction of bright -- and here we mean $r \simeq 24$ as our bright magnitude! -- objects that are stars that would be subject to proper motions, we may simply not be accounting for stellar drift.
The trouble is, of course, we only have \textit{Gaia} to maybe 22nd magnitude, and the HSC-DP1 matches (at least for the moment) don't contain enough objects to overlap with that dynamic range, so we can't cross-check our results.

The conclusion of the HSC-DP1 investigation is, therefore, simply one of patience.
DP2 will provide two orders of magnitude more sky coverage in which to investigate faint LSST astrometry, where we should be able to use the \textit{Gaia} proper motions for bright HSC sources to calibrate out this too-large systematic, if it is driven by epoch baseline differences.
We also do not see the inexplicable power-law slope scaling relation previously observed with OR3, and can either ascribe this to a simulation-specific issue or a pipeline problem that was fixed between the two data reduction runs.

\begin{figure}
  \centering
  \begin{subfigure}[b]{0.95\textwidth}
  \centering
  \includegraphics[width=\columnwidth]{{histogram_mag_vs_sig_vs_snr_object_G_95-25}.pdf}
  \caption{Relationship between magnitude, SNR, and precision for the \texttt{Object} table in the ``95 -25'' pointing.}
  \end{subfigure}
  \hfill
  \begin{subfigure}[b]{0.95\textwidth}
  \centering
  \includegraphics[width=\columnwidth]{{histogram_mag_vs_sig_vs_snr_source_C_47Tuc}.pdf}
  \caption{Comparison relationship between magnitude, SNR, and precision for \texttt{Source} table detections in the ``47 Tuc'' pointing.}
  \end{subfigure}
  \caption{Relationship between magnitude, SNR, and pipeline-derived astrometric precision.
           Left-hand panels show the precision vs log-inverse SNR, the middle panels magnitude vs log-inverse SNR, while the right-hand panels show precision vs magnitude.
           Note in particular the spread in each case; for a uniform set of observations, these three relationships should be completely determined.
           Precision is tightly related to SNR through the image resolution, which, in the case of a ground-based telescope, is set by the atmospheric seeing, with a few examples overplotted on the left-hand panel.}
  \label{fig:mag-snr-unc-relations}
\end{figure}

\subsection{Adding SNR to the \texttt{Object} Table Error-Precision Relation}
\label{sec:object-snr-error-precision}
As mentioned in Section \ref{sec:object-err-precision}, and highlighted in Figure \ref{fig:object_summary}, we see evidence of variation in the magnitude-SNR relationship in at least one of the pointings described above.
This is, ultimately, due to the non-uniformity of the DP1 observations.
This can be seen in figure 6 of \citet{Guy2025}, where the variations in pointing orientation and dither pattern can result in ``five-sigma'' detection limits that vary by a magnitude or more.
Hence, while we step naively through a series of magnitude bins and expect the SNR to step with it, there may have been cases where the detections we were seeing -- particularly in cases such as 47 Tuc where the pipeline failed due to detection overlap -- were systematically in different places on the observations.
These objects would therefore not at the ``expected'' SNR that you would observe for a uniform dataset, such as LSST will eventually produce.
The case of the ``95 -25'' pointing is highlighted in Figure \ref{fig:mag-snr-unc-relations}, contrasted with an example of a much more uniform \texttt{Source} table.
For the \texttt{Object} table detections there is significant spread in the measured magnitude, SNR, and covariance, all of which should be -- as the second result highlights -- very closely related to one another in an ideal set of observations.
This was also one possible suggestion as to the case of the odd behaviour seen in the OR3 \texttt{Object} table, and thus is worth probing further for that reason as well.

We therefore conducted a similar, but slightly changed investigation into the \textit{Gaia}-DP1 errors; we did not repeat the investigation for the \textit{Hubble} truth table equivalent due to small sky overlap.
Here we extracted the \texttt{inputCount} column for our filters -- e.g. \texttt{r\_inputCount} -- in which was stored the number of individual visit images that went into the area of sky in which the detection was made.
We then repeated the entire experiment from Section \ref{sec:object-err-precision} again, but added this coadd depth-equivalent $N$ as an additional dimension to the slices made.

\begin{figure}
    \centering
    \begin{subfigure}[b]{0.45\textwidth}
        \centering
        \includegraphics[width=\textwidth]{{snr_coaddn_r_17A_v2}.pdf}
        \caption{SNR vs coadd count for $r=17$, ``Seagull'' pointing.}
    \end{subfigure}
    \hfill
    \begin{subfigure}[b]{0.45\textwidth}
        \centering
        \includegraphics[width=\textwidth]{{snr_coaddn_r_22.5C_v2}.pdf}
        \caption{SNR vs coadd count for $r=22.5$, ``47 Tuc'' pointing.}
    \end{subfigure}
    \vspace{1em}  % Vertical space between rows
     \begin{subfigure}[b]{0.45\textwidth}
        \centering
        \includegraphics[width=\textwidth]{{snr_coaddn_r_24.5G_v2}.pdf}
        \caption{SNR vs coadd count for $r=24.5$, ``95 -25'' pointing.}
    \end{subfigure}
    \hfill
    \begin{subfigure}[b]{0.45\textwidth}
        \centering
        \includegraphics[width=\textwidth]{{snr_coaddn_r_26D_v2}.pdf}
        \caption{SNR vs coadd count for $r=26$, ``EuDFS'' pointing.}
    \end{subfigure}
     \caption{Scaling relationship between SNR and coadd count.
              2D histograms show the derived SNRs and number of images that contributed to the coadd at the position of a set of detections.
              Red lines show the scaling law $y = p \sqrt{x}$, for the average $p = y / \sqrt{x}$ of all data shown.
              Each sub-figure shows a different magnitude slice (each $r\pm0.1$) and DP1 pointing, as indicated in the sub-figure caption.}
     \label{fig:snr_coaddn_relation}
\end{figure}

Given the further reduction in number statistics this resulted in for DP1, we increased the magnitude slice cut to $r\pm0.25$; we also decided not to make a full least-squares fit of the magnitude-astrometric precision-coadd count slices.
Instead, we compute the error-based scatter metric from the various computable parameters associated with the underlying Rayleigh distribution: the standard deviation of the data; their mean, median, and mode; and a maximum-likelihood estimator (MLE) appropriate to the distribution \citep{siddiqui1964}.
Since we further remove the false-positive rejection with the lack of multi-model fitting, we eliminate objects from the statistic if they fail a colour cut and if they are more than $8\sigma{'}$ apart, computing the initial $\sigma{'}$ from the median separation as an outlier-robust metric.
With these two cuts all five of our metrics are in good agreement -- the mode suffers from binned-data resolution and is noisy but should be insensitive to model mis-specification like the median, in contrast with the standard deviation, mean, and MLE which are sensitive to outliers -- and hence have confidence in the derived error-based scatter scale parameters.
Ultimately having to choose one, we use the median in what follows.

The very first question to confirm is whether the (photometric) SNRs scale in the right way with increasing coadd depth.
Figure \ref{fig:snr_coaddn_relation} shows the relationship between SNR and number of images contributing to the coadd, $N$, for a selection of pointings and magnitude cuts.
Given that signal-to-noise ratios ultimately boil down to counting statistics, each image we add to the first should increase the counts linearly but the uncertainties by the square-root of the new counts, and hence we expect a scaling between SNR and $N$ that goes as $\sqrt{N}$.
The red line in each sub-figure of Figure \ref{fig:snr_coaddn_relation} shows $y = p \sqrt{x}$, with $p$ calculated very simply from the mean of the SNR of each object divided by the square-root of its coadd count.
These show good agreement across all magnitude slices and every pointing, so we can have faith in the photometric SNRs, at least as they scale with number of stacked images that go into their computation.

\begin{figure}
    \centering
    \begin{subfigure}[b]{0.45\textwidth}
        \centering
        \includegraphics[width=\textwidth]{{mag_sigma_coaddn_r_Median_r18_v2}.pdf}
        \caption{Error-based scatter vs coadd count for all DP1 pointings, for $r=17$.}
    \end{subfigure}
    \hfill
    \begin{subfigure}[b]{0.45\textwidth}
        \centering
        \includegraphics[width=\textwidth]{{mag_sigma_coaddn_r_Median_r19.5_v2}.pdf}
        \caption{Error-based scatter vs coadd count for all DP1 pointings, for $r=19.5$.}
    \end{subfigure}
    \vspace{1em}  % Vertical space between rows
     \begin{subfigure}[b]{0.45\textwidth}
        \centering
        \includegraphics[width=\textwidth]{{mag_log_sigma_coaddn_r_Median_r17_v2}.pdf}
        \caption{Error-based scatter vs coadd count for all DP1 pointings, for $r=17$.
                 Note the logarithmic axes, highlighting the low-SNR regime.}
    \end{subfigure}
    \hfill
    \begin{subfigure}[b]{0.45\textwidth}
        \centering
        \includegraphics[width=\textwidth]{{mag_log_sigma_coaddn_r_Median_r20_v2}.pdf}
        \caption{Error-based scatter vs coadd count for all DP1 pointings, for $r=20$.
                 Note the logarithmic axes, highlighting the low-SNR regime.}
    \end{subfigure}
     \caption{Scaling relationship between astrometric DP1-\text{Gaia} cross-match separation distance scatter and coadd depth.
              In each sub-figure, the computed scaling parameter for the underlying Rayleigh distribution is plotted against coadd image number, for each DP1 pointing in different colours, at a fixed magnitude.
              The dash-dot line shows a fit to the relation of $A N^{-c}$, while the dotted line shows a sum-in-quadrature model of the form $\sqrt{(A N^{-0.5})^2 + c^2}$.
              The units of all y-axes are arcseconds.}
     \label{fig:scatter_coaddn_relation}
\end{figure}

The next question is what the expected scaling relation between astrometric centroid precision and coadd image count is.
Assuming that centroid precision is inversely proportional to SNR (see Appendix \ref{sec:hsc_appendix}, but also e.g. \citealp{King:1983aa}), then we have a scaling relation that is inversely proportional to the square-root of $N$, following Figure \ref{fig:snr_coaddn_relation}.
However, the scaling relations between astrometric centroid confidence, in the form of error-based scatter, and coadd image count that we observe in Figure \ref{fig:scatter_coaddn_relation} do not obey a pure $\sqrt{N}$ scaling law.
Each slice in magnitude -- to hold the fluxes and hence counts constant -- shows the relationship between our derived Rayleigh-based $\sigma$ parameter, a function of e.g. the median \textit{Gaia}-DP1 measurement error, and the number of images contributing to the coadds for those objects, for all DP1 pointings.
For each pointing we overplot two potential models: a sum-in-quadrature of the expected inverse-square-root scaling with an unknown systematic uncertainty (dotted lines), and a relaxed power-law fit (dash-dot lines).
Given the low number of objects in each bin and hence the noise in each individual fit we cannot distinguish between these two models, but both fit the data reasonably well; certainly better than the fixed inverse-square-root law.

\begin{figure}
  \centering
  \includegraphics[width=0.95\columnwidth]{{mag_scaling_fits_r_Median}.pdf}
  \caption{Summary statistics of the magnitude-precision-SNR error-based scatter model.
           The top two panels show the values we can derive from the sum-in-quadrature model $\sqrt{(A N^{-0.5})^2 + c^2}$, with the $N=1$ case of $\sqrt{A^2 + c^2}$ on the left and the $N \to \infty$ $c$-only limit on the right.
           The bottom two panels show values from the power-law model $A N^{-c}$, again showing the case for $N=1$ on the left and the exponent $c$ on the right.
           Each is plotted as a function of magnitude.}
  \label{fig:scatter_coaddn_parameters}
\end{figure}

From these models we can extract three values.
First, in both cases we can calculate the expected error-based scatter metric for $N=1$, i.e., a case where no coaddition has been performed at all.
For the sum-in-quadrature model we can also compute the case of $N \to \infty$, i.e., the ``plateau'' value of the systematic precision.
In the case of the power-law model we can also report on the exponent itself.
These values are shown in Figure \ref{fig:scatter_coaddn_parameters}.

The two $N=1$ model-based precisions are in good agreement.
Each shows a variation from 5 mas, at $16 < r < 17$, to around 10-15 mas, both brighter -- above the saturation limit of the LSST -- than 16th magnitude and fainter than 17th.
We see a similar shape to the best-fit missing systematic precisions, varying from 2 mas up towards 10 mas.
On the other hand, the equally valid power-law slope varys from relationships of the form $\sigma \propto N^{-0.25}$ at $r\simeq16$ towards flatter and flatter slopes at faint magnitudes, perhaps as low as $\sigma \propto N^{-0.1}$ by 21st magnitude.
The values brighter than 16th magnitude here are questionable due to the negative power-law slopes, an artefact of some noisy data not being handled as robustly by this model than by the sum-in-quadrature model; we didn't set a positive-$c$ constraint but do not try to interpret such unphysical results either.

This then goes someway to explaining the point made previously about the disconnect between SNR and magnitude.
If, at each magnitude slice used in e.g. the ``95 -25'' pointing analysis (cf. e.g. Figures \ref{fig:excdfs_obj_h_24.4} and \ref{fig:object_summary}), we are consistently selecting from high-coadd-depth regions of the observed pointing, then we may be biased towards detections that have overestimated precisions due to their not following the $N^{-0.5}$ law for whatever reason.
Alternatively, if the sum-in-quadrature model is right, then we may have a ``missing systematic'' precision, or scatter in measured positions, that is variant with brightness; this would explain the fact that we have error-based scatters well above quoted precisions for objects too faint to merely be a traditional, roughly few mas, systematic floor from things like World Coordinate System (WCS) uncertainty, independent of object flux.

\subsection{The \texttt{Source} Table Error-Precision Relation}
Although Section \ref{sec:object-snr-error-precision} touched upon the $N=1$, \texttt{Source} table-regime within the \texttt{Object} table, we can, of course, characterise the visit-level catalogue position-scatter vs precision relationship directly.
This catalogue is the time-series component of the LSST, with each exposure having its detections extracted and converted to positions and brightnesses that can be traced through time.
Here we therefore have many times more datapoints than the case of the \texttt{Object} table, getting some way towards a set of repeat measurements for each real sky source for which we can compare precisions and errors.

The experiment follows similarly to the previous details laid out in Sections \ref{sec:validating-positions-precisions} and \ref{sec:object-err-precision}.
From 14th through 24th magnitude -- in both cases of \textit{Gaia} and HSC as our truth table from which to derive measurement errors, since LSST now has the same completeness limit -- we take slices in magnitude, in both increments of 0.1 magnitudes and bin widths of 0.1 magnitudes again.
Cuts are made to ensure outlier rejection of pipeline-derived precisions, and slices are removed for low-number statistics, although since we have many times the number of unique detections across all visits this happens much less frequently.
The only minor difference is that we use the $g$ band detections rather than the $r$ band this time, not being so concerned with the coverage as we were in the \texttt{Object} table analysis, preferring the bluer filter to minimise colour effects where possible.

\begin{figure}
    \centering
    \begin{subfigure}[b]{0.45\textwidth}
        \centering
        \includegraphics[width=\textwidth]{{auf_fits_A_Seagull_source_gaia_g16.6}.pdf}
        \caption{DP1-\textit{Gaia} error-based scatter distributions and models of precisions for the ``Seagull'' pointing and $g=16.6$.}
        \label{fig:seagull_source_g_16.6}
    \end{subfigure}
    \hfill
    \begin{subfigure}[b]{0.45\textwidth}
        \centering
        \includegraphics[width=\textwidth]{{auf_fits_D_EuDFS_source_gaia_g21.2}.pdf}
        \caption{DP1-\textit{Gaia} error-based scatter distributions and models of precisions for the ``EuDFS'' pointing and $g=21.2$.}
        \label{fig:eudfs_source_g_21.2}
    \end{subfigure}
    \vspace{1em}  % Vertical space between rows
     \begin{subfigure}[b]{0.45\textwidth}
        \centering
        \includegraphics[width=\textwidth]{{auf_fits_A_Seagull_source_gaia_g22.3}.pdf}
        \caption{DP1-\textit{Gaia} error-based scatter distributions and models of precisions for the ``Seagull'' pointing and $g=22.3$.}
        \label{fig:seagull_source_g_22.3}
    \end{subfigure}
    \hfill
    \begin{subfigure}[b]{0.45\textwidth}
        \centering
        \includegraphics[width=\textwidth]{{auf_fits_F_ExCDFS_source_hubble_g24.0}.pdf}
        \caption{DP1-HSC error-based scatter distributions and models of precisions for the ``ExCDFS'' pointing and $g=24$.
                 Note the slightly elongated non-Gaussian tail to the distribution, possibly due to missing proper motion correction.}
        \label{fig:excdfs_source_h_24}
    \end{subfigure}
     \caption{Comparison between astrometric precision model and distribution of measured-truth position separations.
              All symbols and lines have the same meaning as Figure \ref{fig:95-25_obj_g_18.9}. 
              Each sub-figure's caption describes the particular magnitude slice, pointing, and ``truth'' table used.}
     \label{fig:source-table-error-histograms}
\end{figure}

\begin{figure}
  \centering
  \includegraphics[width=0.95\columnwidth]{{hubble_corrected_lsst_errors_source}.pdf}
  \caption{Summary of pipeline-derived uncertainties and experimental error-based scatter for DP1 astrometry.
           All colours and symbols have the same meaning as in Figure \ref{fig:object_summary}.
           The HSC data have had the non-negligible \textit{Hubble} astrometric precision removed while we assume the \textit{Gaia} precisions effectively do not contribute to the nearest neighbour distances.
           The left-hand panel shows the data in a linear scale while the right-hand panel shows the same datapoints on a logarithmic scale, to highlight the very brightest observations.}
  \label{fig:source_summary}
\end{figure}

We show a few example magnitude-pointing combinations in Figure \ref{fig:source-table-error-histograms}.
We see a very similar story to the \texttt{Object} table analysis, cf. Figures \ref{fig:95-25_obj_g_18.9} and \ref{fig:excdfs_obj_h_24.4}.
At bright magnitudes (Figure \ref{fig:seagull_source_g_16.6}), we see the typical underestimated pipeline-derived precision when compared to the measurement error distribution.
At intermediate brightnesses, however, we see a mix of results.
We have at least one pointing (``Seagull'', Figure \ref{fig:seagull_source_g_22.3}) in which, down below 22nd magnitude, we have agreement between pipeline precision and measurement scatter.
However, all others show varying degrees of the story told by Figure \ref{fig:eudfs_source_g_21.2}, in which we continue to see the precisions be lower than the standard deviation of the errors.
Finally, the single \textit{Hubble}-based pointing we have good overlap with in the DP1 dataset also shows the now-typical story of precision underestimation, even when the HSC astrometric precisions are subtracted from the error distributions (as discussed in Section \ref{sec:object-err-precision} and Appendix \ref{sec:hsc_appendix}).
These results are summarised in Figure \ref{fig:source_summary}.

\begin{figure}
  \centering
  \includegraphics[width=0.95\columnwidth]{{summary_mn_sky_source_gaia}.pdf}
  \caption{Parameterised model of DP1 \texttt{Source} table uncertainties, as determined by DP1-\textit{Gaia} error measurements for all DP1 pointings.
           The top panels show the hyper-parameters $m$ and $n$ respectively from the model $\sigma_\mathrm{err}^2 = (m\sigma_\mathrm{covar})^2 + n^2$.
           Each pointing is represented on an equatorial sky projection, with Galactic latitude $b = 0$ and $\lvert b \lvert = 20$ shown in the solid lines and Ecliptic latitudes of zero and $\pm20$ shown in the dotted lines.
           The bottom two plots show $m$ and $n$ respectively in the left and right panels again, but this time plotted as a function of average DP1 density.}
  \label{fig:source_mn}
\end{figure}

It is clearer from Figure \ref{fig:source_summary} that each pointing has a different scaling relation between position-based scatter and its pipeline-derived precisions, with the black plusses representing the case of the Seagull Nebula highlighted in Figures \ref{fig:seagull_source_g_16.6} and \ref{fig:seagull_source_g_22.3}.
This contrasts with the results \citet{sitcomtn-159} found for the simulated LSST data of consistently well-behaved error-precision relationships that only required a systematic uncertainty to be included at bright fluxes (e.g. \citealp{sitcomtn-159}, figure 5).
Following the same investigation steps, we can also reproduce the hyper-parameter fits (as were previously also discussed in Section \ref{sec:object-err-precision}).
These are shown in Figure \ref{fig:source_mn}.

The linear slopes seen in Figure \ref{fig:source_summary} become clear when fit with a model of $\sigma_\mathrm{err}^2 = (m\sigma_\mathrm{covar})^2 + n^2$, with $m$ in the range approximately 1.2-1.9.
We see necessary systematic precisions (cf. Figure \ref{fig:seagull_source_g_16.6}) of roughly 8-22 mas, larger than the 3-7 mas quoted by \citet{sitcomtn-159} for simulated LSST data.
However, we do see a negative trend between $n$ and typical LSST field density\footnote{Note that several of the pointings, in particular that centered on 47 Tuc, suffer significant issues with source extraction in dense regions, and hence are systematically below their expected density, but we make no correction for this.}, as seen in the lower-right panel of Figure \ref{fig:source_mn}.
Here we see a tentative explanation for larger systematic precisions than previously quoted: the lowest densities probed by OR3 were 20,000 sources per square degree, while all but two of our quoted densities are below that value.
In essence, there \textit{may} be a continuous function that connects the systematic precisions derived by both the OR3 and DP1 datasets, but we need to probe higher-density regions -- and successfully extract sources at these higher densities -- to be able to confirm.

On the other hand, $m$ shows a very weak negative trend with density, in a reversal of that reported by \citet{sitcomtn-159} for OR3.
However, as the two datasets probe very different density regimes and used very different versions of the LSST Science Pipelines, we cannot say for sure whether this is a change driven by software or if it is a real phenomenon.
We, again, require a single complete dataset that spans from sub-10,000 sources per square degree to the 100,000s of sources per square degree previously detected to know for sure.

Our single HSC pointing overlap further corroborates the under-estimation of LSST precisions when shown in full in Figure \ref{fig:source_summary}.
As can be seen, the two analyses of DP1 \texttt{Source} table -- using \textit{Gaia} and the HSC as our input ``truth'' respectively -- agree in the area of overlap in dynamic range we observe, before \textit{Gaia}'s completeness limit causes a drop off in counterpart associations.
We see a similar result of $m \simeq 1.2$ for a hyper-parameter fit, although this comes with an $n \approx 60$ mas systematic, and should therefore be viewed with some skepticism.
All of the arguments made in Section \ref{sec:object-err-precision} regarding the quality of the HSC data at bright magnitudes still apply.

\section{Conclusions, and Towards Data Preview 2}
The main conclusion from the work presented here is that there is evidence that the LSST Science Pipelines consistently underestimate the precision with which any individual detection -- object or source -- is centroided in the images generated by the Rubin Observatory as part of what will very soon be the LSST.
We report varying degrees of systematic precision that must be combined with the reported uncertainties to account for the scatter in measured LSST positions, as measured against both \textit{Gaia} and the \textit{Hubble} Source Catalog, at bright magnitudes.
However, this trend continues to faint magnitudes, where we would expect photon-counting statistics to dominate.
This suuggests that perhaps some component of the LSST Science Pipelines is failing to take into account a crucial aspect of measurement-belief quantification, or maybe that there are additional forces at play, either within the telescope or the software and algorithms, leading to deviations away from ``true'' that need to be corrected for.
An investigation into the trend of astrometric position scatter vs pipeline-derived precisions as a function of coadd depth gives some evidence for an a-typical scaling relation, providing one possible explanation as to the too-large \texttt{Object} table precisions, although these issues are seen in both \texttt{Object} and \texttt{Source} tables.
These too-precise uncertainties, or too-large measurement-error scatters equivalently, have significant knock-on effects for a wide-ranging number of downstream science cases, such as proper motion and parallax fits for stellar drift, progenitor determination for transient events, and the broad sub-field of catalogue cross-matching and counterpart association.

These results are, however, preliminary in every sense.
The LSST Science Pipelines have progressed since these data were processed, and include more sophisticated modelling of a number of contributions to detected-position deviations that should reduce the lengthscales over which the measurement errors are distributed, bringing error and precision into closer alignment immediately.
And, of course, Data Preview 1 is just that: the first preview of what is to come.
Data Preview 2 will contain at least two orders of magnitude more sky coverage, to similar if not deeper co-added completeness limits, across a more homogeneous region of the sky.
In addition, it will benefit from the larger field-of-view of the main LSSTCam, as opposed to the LSSTComCam used for the observations that make up DP1, which will greatly reduce any systematic effects that might be related to the size of the image, such as the number of detections used to generate the WCS solution.
These data will be crucial for determining the most up-to-date accuracy and precision of the LSST Science Pipelines' measured astrometric positions and covariances.
They will provide early characterisation of the systematics and additional un-modelled sources of astrometric measurement uncertainty that will be needed for the first major LSST Data Release, due in 2028.

\appendix

\section{Acknowledgements}

This material is based upon work supported in part by the National Science Foundation through Cooperative Agreements AST-1258333 and AST-2241526 and Cooperative Support Agreements AST-1202910 and AST-2211468 managed by the Association of Universities for Research in Astronomy (AURA), and the Department of Energy under Contract No.\ DE-AC02-76SF00515 with the SLAC National Accelerator Laboratory managed by Stanford University.
Additional Rubin Observatory funding comes from private donations, grants to universities, and in-kind support from LSST-DA Institutional Members.

This work has made use of data from the European Space Agency (ESA) mission \textit{Gaia} (\url{http://www.cosmos.esa.int/gaia}), processed by the \textit{Gaia} Data Processing and Analysis Consortium (DPAC, \url{http://www.cosmos.esa.int/web/gaia/dpac/consortium}). Funding for the DPAC has been provided by national institutions, in particular the institutions participating in the \textit{Gaia} Multilateral Agreement.
Based on observations made with the NASA/ESA Hubble Space Telescope, and obtained from the Hubble Legacy Archive, which is a collaboration between the Space Telescope Science Institute (STScI/NASA), the Space Telescope European Coordinating Facility (ST-ECF/ESAC/ESA) and the Canadian Astronomy Data Centre (CADC/NRC/CSA).

This paper makes use of LSST Science Pipelines software developed by the Vera C. Rubin Observatory. We thank the Rubin Observatory for making their code available as free software at https://pipelines.lsst.io \citep{NDVRO2025,Jenness2022}.
It also uses the \texttt{numpy} \citep{Harris2020}, \texttt{astropy} \citep{2013A&A...558A..33A,2018AJ....156..123A}, \texttt{scipy} \citep{scipy}, \texttt{matplotlib} \citep{matplotlib}, and \texttt{pandas} \citep{mckinney-proc-scipy-2010-pandas} softwares.

\section{On the Astrometric Precision of the \textit{Hubble} Source Catalog, and the Uncertainty of an Aperture Centroid}
\label{sec:hsc_appendix}
In Section \ref{sec:object-err-precision} we consider the joint positional errors of the \textit{Hubble} Source Catalog and LSST's DP1.
However, the HSC data do not contain astrometric precisions, so we cannot confirm directly whether, as in the case of \textit{Gaia}-LSST match separations, our errors are dominated by LSST or not.
Thus, we investigated the scatter between \textit{Gaia} and HSC symmetric, nearest-neighbour matches, as the three-catalogue tiebreak.
Here, for both pointings that have common LSST DP1-HSC sky overlap, we can derive astrometric errors for HSC -- at least brighter than 22nd magnitude in \textit{Gaia}.

Across the seven regions of sky the DP1 pointings cover, we see a consistent scatter between \textit{Gaia} and HSC observations, once the time baseline between the two datasets is accounted for.
In all cases, the standard deviation of the errors\footnote{This refers to the ``$\sigma$'' value of the underlying Gaussian or Rayleigh distribution. We can compute the standard deviation of the radial match separations, but -- while related to the Gaussian or Rayleigh $\sigma$ -- these are not quite identical. As we discuss in Section \ref{sec:object-snr-error-precision}, you can compute the underlying scale parameter $\sigma$ in numerous different ways: the standard deviation, the mean, the median, the mode, or a maximum-likelihood estimator. As all of these give the same answer here, we are confident that the parameter is robust and consistent.} is between 12 and 20 milli-arcseconds.
Given the meta-catalogue generation of the HSC, and \textit{Hubble}'s weaknesses in deriving absolute astrometric solutions, the next question to ask is whether this is a systematic uncertainty or noise-based scatter.
In the first instance, these bright-source errors would extend trivially to fluxes below the \textit{Gaia} completeness limit as they would be independent of the brightness of the objects, while the scatter in counterpart distances being related to the SNR of the objects would require a more complex description.
To determine this, we now consider the mathematical form of the position determined in the HSC, and model its determination precision.

The catalogues use the \texttt{SExtractor} software \citep{Bertin1996} to derive their parameters, which computes centre-of-mass centroids,
\begin{equation}
  \hat{x} = \frac{\sum_i p_i x_i}{\sum_i p_i}.
\end{equation}
In the notation of \citet{King:1983aa}, $p_i = l f_i + \alpha (b - \hat{b})$, with $l$ the total counts of the source, $f_i$ the fraction of the point-spread function (PSF) in the $i$th pixel, $\alpha$ the area of a pixel and $b$ the background counts per unit area; $\hat{b}$ is an estimation of the background flux.
The variance of the aperture-centroid is given by
\begin{equation}
  \sigma_{\hat{x}}^2 = \sum_i \left(\frac{\partial \hat{x}}{\partial p_i}\right)^2 \sigma_{p_i}^2,
\end{equation}
where crucially the variance of the pixel must be the non-background-subtracted variance, so $\sigma_{p_i}^2 = l f_i + \alpha b$.

The differential of $\hat{x}$ with respect to any individual pixel $p_i$ is given by
\begin{align}
\begin{split}
  \frac{\partial \hat{x}}{\partial p_i} &= \frac{\frac{\partial}{\partial p_i} \left(\sum_j p_j x_j\right)}{\sum_j p_j} + \sum_j p_j x_j \frac{\partial}{\partial p_i}\left(\left(\sum_j p_j\right)^{-1}\right)
                                        = \frac{x_i}{\sum_j p_j} - \frac{\sum_j p_j x_j} {\left(\sum_j p_j\right)^{2}}.
\end{split}
\end{align}
Again following the logic of \citet{King:1983aa}, we can consider the limit where pixels become small relative to the PSF, and our sums become integrals:
\begin{equation}
  \sum_i l f_i + \alpha b = \sum_i l \phi_i \alpha + \alpha b = \int\! (l \phi + b)\, \mathrm{d}A,
\end{equation}
with $\phi$ the mathematical representation of the PSF, now with dimensions of inverse-area, and $\int\! \phi\,\mathrm{d}A = 1$, as we assume $\phi$ is normalised.
It should now be more obvious that the second term of the differential of $\hat{x}$ with respect to $p_i$ vanishes, as we can state it as
\begin{equation}
  \sum_i p_i x_i = \sum_i (l f_i + \alpha b) x_i = \sum_i (l \phi_i \alpha + \alpha b) x_i = \int\! (l \phi + b) x\, \mathrm{d}A,
\end{equation}
where, regardless of the exact location of the object (at some arbitrary $x_0$), the implicit integral limits of $\pm\infty$ -- or just any equal, but opposite sign, limits in $x$ around $x_0$ -- mean that by symmetry the entire integral evaluates to zero.

The variance in the centroid position is therefore
\begin{align}
\begin{split}
  \sigma_{\hat{x}}^2 &= \sum_i \left(\frac{\partial \hat{x}}{\partial p_i}\right)^2 \sigma_{p_i}^2 = \sum_i \left(\frac{x_i}{\sum_j p_j}\right)^2 \sigma_{p_i}^2 
                     = \frac{\int\!x^2(l\phi + b)\,\mathrm{d}A}{\left(\int\!(l\phi + (b - \hat{b}))\,\mathrm{d}A\right)^2}.
\end{split}
\end{align}
If we make the assumption that $b = \hat{b}$, and the background is estimated accurately, this simplifies to
\begin{align}
\begin{split}
  \sigma_{\hat{x}}^2 &= \frac{\int\!x^2(l\phi + b)\,\mathrm{d}A}{\left(\int\!l\phi\,\mathrm{d}A\right)^2}.
\end{split}
\end{align}
We can now consider two extremes of the brightness of the central object: either it is so bright that the noise is entirely dominated by the photon counts of the object, or it is very faint compared with the sky background and all pixels have the same weight.

In the first case, we can drop the $b$ term from the numerator, as $b \to 0$.
Assuming the PSF is approximately a Gaussian, as is a reasonable approximation for seeing-limited observations, the integral $\int\!l x^2 \phi\,\mathrm{d}A$ evaluates to $l \sigma^2$, with $\sigma$ related to the full-width at half-maximum (FWHM), or astronomical seeing, of the observations by $\sigma = \mathrm{FWHM} / (2 \sqrt{2 \ln(2)})$.
The denominator trivially evaluates to $l^2$, since $\phi$ is normalised, and hence we obtain
\begin{align}
\begin{split}
  \sigma_{\hat{x}}^2 &= \frac{\int\!x^2l\phi\,\mathrm{d}A}{\left(\int\!l\phi\,\mathrm{d}A\right)^2} = \frac{l \sigma^2}{l^2} = \frac{\sigma^2 \sigma_l^2}{l^2},
\end{split}
\end{align}
where the final equality makes use of the fact that, in the bright-object limit, the variance of the flux is given by
\begin{equation}
  \sigma_l^2 = \sum_i \sigma_{p_i}^2 = \sum_i l f_i = \int\! l \phi\,\mathrm{d}A = l.
\end{equation}
This result is very similar to that derived by \citet{King:1983aa} in the case of least-squares fitting, which has an intuitive explanation in the fact that each pixel -- or infinitesimal unit area, if you make the swap from a sum to an integral -- is implicitly weighted by its own flux in the original equation for the centroid.

For the background-dominated case, we see that $b \gg l$, and we have
\begin{align}
\begin{split}
  \sigma_{\hat{x}}^2 &= \frac{\int\!x^2\,b\,\mathrm{d}A}{\left(\int\!l\phi\,\mathrm{d}A\right)^2},
\end{split}
\end{align}
which is difficult since the numerator contains an integral that does not converge.
If we take an integral in both numerator and denominator over some circle of radius $r$, then we can see
\begin{equation}
  \int_{r{'} < r}\!\phi\,\mathrm{d}A = 1 - \exp\left(-\frac{1}{2} \frac{r^2}{\sigma^2}\right)
\end{equation}
for a Gaussian-approximated PSF, while
\begin{equation}
  \int_{r{'} < r}\! x^2\,\mathrm{d}A = \int_0^{2\pi}\!\int_0^{r} r{'}^3 \cos^2{\theta}\,\mathrm{d}r{'}\,\mathrm{d}\theta = \pi \frac{r^4}{4}.
\end{equation}
This results in
\begin{align}
\begin{split}
  \sigma_{\hat{x}}^2 &= \frac{\int_{r{'} < r}\! x^2b\,\mathrm{d}A}{\left(\int_{r{'} < r}\!l\phi\,\mathrm{d}A\right)^2} = \frac{b \pi r^4}{4 l^2 \left(1 - \exp\left(-\frac{1}{2} \frac{r^2}{\sigma^2}\right)\right)^2} = \frac{r^2 \sigma_l^2}{4l^2\left(1 - \exp\left(-\frac{1}{2} \frac{r^2}{\sigma^2}\right)\right)^2},
\end{split}
\end{align}
where again we have used the variance of the flux in the background-dominated regime,
\begin{equation}
  \sigma_l^2 = \sum_i \sigma_{p_i}^2 = \sum_i \alpha b= \int_{r{'} < r}\! b\,\mathrm{d}A = b \pi r^2.
\end{equation}
This result, interestingly, is very similar to that given by \citet{Naylor:1998aa} through different considerations and only focusing on the SNR of the photometric measurement.

For the case of \textit{Hubble}, for example in its ACS F814W filter, the PSF FWHM is of order 0.12 arcseconds, while the aperture radius used to derive its photometry was 0.15 arcseconds.
Thus, in the first, bright-source case, we have a case of $\sigma_{\hat{x}} \approx 0.05\,\mathrm{arcsec} \times \mathrm{SNR}^{-1}$, while in the background-dominated case we have $\sigma_{\hat{x}} \approx 0.075\,\mathrm{arcsec} \times \mathrm{SNR}^{-1}$.
This shows a similar indifference to brightness regime that \citet{King:1983aa} demonstrates for the least-squares case, but more crucially provides an order-of-magnitude estimate for pure centroid constraint precisions.

Given the brightness of the \textit{Hubble} objects that we can match to \textit{Gaia}, we must have SNRs that are large -- at least of order 50, corroborated by MADs of order 0.02 magnitudes -- and hence our centroid precisions would be of order 1 mas.
As we see \textit{Gaia}-\textit{Hubble} error deviations of order 10 mas, we must therefore conclude that these errors are driven by systematics beyond those of individual noise-based deviations.
We therefore apply a simple, systematic-uncertainty-based centroid uncertainty to our DP1-HSC errors to approximate the LSST component of the scatter, subtracting in quadrature the \textit{Gaia}-HSC errors determined for bright \textit{Hubble} observations.
Hence, taking the average of our various determinations of the scatter in \textit{Gaia}-HSC errors, we assume a constant 15 mas precision on the \textit{Hubble} data.


% Include all the relevant bib files.
% https://lsst-texmf.lsst.io/lsstdoc.html#bibliographies
\section{References} \label{sec:bib}
\renewcommand{\refname}{} % Suppress default Bibliography section
\bibliography{local,lsst,lsst-dm,refs_ads,refs,books}

% Make sure lsst-texmf/bin/generateAcronyms.py is in your path
\section{Acronyms} \label{sec:acronyms}
\input{acronyms.tex}
% If you want glossary uncomment below -- comment out the two lines above
%\printglossaries





\end{document}
